\section{Conclus\~{a}o e Trabalhos Futuros}
\label{sec:conclusao}

\subsection{S\'{i}ntese das Contribui\c{c}\~{o}es}

Este artigo apresentou o \textbf{Sentinela Clandestino}, uma arquitetura de intelig\^{e}ncia artificial na borda voltada \`{a} conten\c{c}\~{a}o em tempo real de ecossistemas digitais de apostas clandestinas no Brasil. Confrontando uma crise que drena entre R\$~6 e R\$~10 bilh\~{o}es mensais da renda das fam\'{i}lias para esquemas criminosos desregulamentados, o trabalho demonstrou que as defesas convencionais falham sistematicamente diante do dinamismo de algoritmos DGA, cobran\c{c}as Pix rotativas e f\'{a}bricas automatizadas de malware polim\'{o}rfico (MAMF).

A solu\c{c}\~{a}o proposta estabelece tr\^{e}s avan\c{c}os t\'{e}cnicos fundamentais:
\begin{enumerate}
    \item \textbf{Fus\~{a}o Sem\^{a}ntico-Topol\'{o}gica na Borda:} Integra a interpreta\c{c}\~{a}o sem\^{a}ntica de termos contratuais abusivos e manifestos de aplica\c{c}\~{o}es com a an\'{a}lise estrutural de resumos de grafos de DNS passivo (pDNS), preenchendo a lacuna existente entre filtros de rede est\'{a}ticos e sistemas de modera\c{c}\~{a}o em nuvem;
    \item \textbf{Pipeline H\'{i}brido de Dois Est\'{a}gios com Privacidade Absoluta:} A combina\c{c}\~{a}o de um modelo de decis\~{a}o Sistema~1 de pesos abertos (\textit{LAYA}) operando em $< 35\text{ms}$ com um SLM quantizado em INT4 acionado sob demanda garante que a experi\^{e}ncia de navega\c{c}\~{a}o n\~{a}o sofra degrada\c{c}\~{a}o percept\'{i}vel, assegurando que $100\%$ das infer\^{e}ncias ocorram dentro do aparelho, sem envio de telemetria para a nuvem ($\mathcal{X}_{\text{telemetria}} = \emptyset$);
    \item \textbf{Canaliza\c{c}\~{a}o do Interesse Pr\'{o}prio do Apostador:} Ao adotar a l\'{o}gica comportamental de painel de evid\^{e}ncias de risco (paradigma Waze), o sistema neutraliza a resist\^{e}ncia psicol\'{o}gica \`{a} censura paternalista, municiando a v\'{i}tima com provas imediatas de calote de saque e presen\c{c}a de \textit{trojans} banc\'{a}rios.
\end{enumerate}

\subsection{Trabalhos Futuros}

Como desdobramentos cient\'{i}ficos e tecnol\'{o}gicos desta investiga\c{c}\~{a}o, planejam-se tr\^{e}s linhas centrais de pesquisa:
\begin{itemize}
    \item \textbf{Sincroniza\c{c}\~{a}o Privada via Aprendizado Federado:} Desenvolvimento de protocolos de Aprendizado Federado (\textit{Federated Learning}) com Privacidade Diferencial para agregar atualiza\c{c}\~{o}es de filtros de Bloom de novos espelhos de amea\c{c}as diretamente entre dispositivos na borda, sem expor o perfil de navega\c{c}\~{a}o individual de nenhum cidad\~{a}o a servidores centrais;
    \item \textbf{Adapta\c{c}\~{a}o Cont\'{i}nua a Dialetos e G\'{i}rias Regionais:} Refinamento de conjuntos de dados e fine-tuning do SLM para capturar a evolu\c{c}\~{a}o dialetal das express\~{o}es de aliciamento de apostas nas diferentes regi\~{o}es brasileiras (e.g., ``plataforma pagando no bug'', ``alavancagem de banca'', ``sala vip de sinais'');
    \item \textbf{Piloto Experimental e Estudos de Campo:} Condu\c{c}\~{a}o de ensaios de usabilidade e conten\c{c}\~{a}o com grupos focais de usu\'{a}rios volunt\'{a}rios, em parceria com \'{o}rg\~{a}os p\'{u}blicos de prote\c{c}\~{a}o ao consumidor (Procon) e institutos de pesquisa em seguran\c{c}a cibern\'{e}tica para o bem social (\textit{AI for Good}).
\end{itemize}